\scanpage{8}
\mspage{-2.33-}

then:

\[
\begin{aligned}
\handwritten{\unsure{\nu|v}}^{M\handwritten{\otimes} f^*L^n}(x,\handwritten{\delta})&=\handwritten{\nu}^{M\handwritten{\otimes}f^*L^{\handwritten{n_0}}}(x,\handwritten{\delta})+\handwritten{\nu}^{f^*L^{\handwritten{n-n_0}}}(x,\handwritten{\delta})\\
&=\handwritten{\nu}^{M\handwritten{\otimes}f^*L^{\handwritten{n_0}}}(x,\handwritten{\delta})+(n-n_0)\handwritten{\nu}^{L}(f(x),\handwritten{\delta}).
\end{aligned}
\]

First of all, apply the first part of Proposition 2.17 to $\handwritten{\nu}^{M\handwritten{\otimes}f^*L^{\handwritten{n_0}}}$. We conclude that, for some $K$:

\[
\handwritten{\nu}^{M\handwritten{\otimes}f^*L^{\handwritten{n_0}}}(x,\handwritten{\delta})\geq-K.
\]

Secondly, apply the second part of Proposition 2.17 to $\handwritten{\nu}^L(y,\handwritten{\delta})$, for $y\in S=Y^s_{(\unsure{o|0})}(L)_k$. By\unsure{$\shortparallel$|n} Theorem 2.1, all these numbers are strictly positive. Therefore, by the Proposition:

\[
\handwritten{\nu}^L(y,\handwritten{\delta})\geq\epsilon>0
\]

for all $y\in Y^s_{(o)}(L)_k$, and $\handwritten{\delta}\in\Delta(G)$. Therefore, if $n\geq\handwritten{\left[}n_0+\frac{K}{\epsilon}+1\handwritten{\right]}$, we conclude

\[
\handwritten{\nu}^{M\handwritten{\otimes}f^*L^n}(x,\handwritten{\delta})>0
\]

whenever $x\in f^{-1}\handwritten{\{}Y^s_{(o)}(L)\handwritten{\}}_k$. Therefore, by Theorem 2.1, all such $x$ are in $X^{s\unsure{{}!|{}'|{}}}_{(o)}(M\handwritten{\otimes} f^*L^n)$. \uline{QED}

% archive sheet number: 7
